\documentclass[10pt]{article}
\usepackage[T1]{fontenc}
\usepackage[utf8]{inputenc}
\usepackage[french]{babel}
\usepackage{amsmath,amssymb,mathrsfs,mathpazo}
\usepackage{geometry}
\geometry{a4paper,margin=2cm}
\usepackage{graphicx,array,tabularx,tabularray,colortbl,multirow,multicol,booktabs,enumitem,pifont,ulem,fancybox,diagbox,cancel}
\usepackage{pstricks,pstricks-add,pst-plot,pst-tree,pst-node,pst-text,pst-func,pst-math,pst-bspline,pst-3dplot,pst-eucl,pst-solides3d}
\let\pscalculate\undefined
\usepackage{pst-bezier}
\usepackage{tikz,pgfplots,tkz-tab}
\usetikzlibrary{arrows,arrows.meta,patterns,calc,positioning,trees,decorations.markings,shapes.geometric,shapes.misc,shapes.arrows,intersections,decorations.pathmorphing,decorations.pathreplacing,backgrounds,shadows,babel}
\pgfplotsset{compat=1.18}
\usepgfplotslibrary{fillbetween}
\usepackage[np]{numprint}
\usepackage{listings,makecell,pas-tableur,twoopt}
\let\barre\undefined
\usepackage[pstricks]{tablvar}
\definecolor{darkgreen}{RGB}{0,100,0}
\providecommand{\red}{\color{red}}
\providecommand{\blue}{\color{blue}}
\providecommand{\green}{\color{green}}
\providecommand{\black}{\color{black}}
\providecommand{\white}{\color{white}}
\providecommand{\gray}{\color{gray}}
\providecommand{\cyan}{\color{cyan}}
\providecommand{\magenta}{\color{magenta}}
\providecommand{\yellow}{\color{yellow}}
\providecommand{\violet}{\color{violet}}
\providecommand{\orange}{\color{orange}}
\providecommand{\euro}{\texteuro}
\providecommand{\EUR}{\texteuro}
\providecommand{\ieme}{\textsuperscript{e}}
\providecommand{\degres}{\ensuremath{^\circ}}
\providecommand{\degree}{\ensuremath{^\circ}}
\providecommand{\vv}[1]{\overrightarrow{#1}}
\providecommand{\Coord}[1]{#1}
\providecommand{\decosix}{}
\providecommand{\starredbullet}{\ensuremath{\bullet}}
\pagestyle{empty}
\setlength{\parindent}{0pt}
\usepackage{adjustbox}\DeclareUnicodeCharacter{2236}{:}\DeclareUnicodeCharacter{2190}{\ensuremath{\leftarrow}}
\providecommand{\R}{\mathbb{R}}
\providecommand{\N}{\mathbb{N}}
\providecommand{\Z}{\mathbb{Z}}
\providecommand{\Q}{\mathbb{Q}}
\providecommand{\D}{\mathbb{D}}
\providecommand{\C}{\mathbb{C}}
\providecommand{\e}{{\mathrm{e}}}
\providecommand{\vect}[1]{\overrightarrow{#1}}
\providecommand{\vectt}[1]{\overrightarrow{\mathrm{#1}}}
\providecommand{\vvt}[1]{\overrightarrow{\mathrm{#1}}}
\providecommand{\np}[2][]{{#2~#1}}
\providecommand{\vv}[1]{{\overrightarrow{#1}}}
\renewcommand{\d}{{\mathrm{d}}}
\providecommand{\degres}{{^\circ}}
\providecommand{\degre}{{^\circ}}
\providecommand{\blue}{\color{blue}}
\providecommand{\red}{\color{red}}
\providecommand{\enskip}{\quad}
\providecommand{\llbracket}{\lbrack\!\lbrack}
\providecommand{\rrbracket}{\rbrack\!\rbrack}
\providecommand{\euro}{{\text{€}}}
\providecommand{\no}{{\text{n°}}}
\providecommand{\interval}[2]{{[#1;#2]}}
\providecommand{\intervFO}[2]{{[#1;#2[}}
\providecommand{\dotfill}{\ldots}

\providecommand{\ds}{\displaystyle}
\providecommand{\textding}[1]{#1}
\providecommand{\Oij}{\(\mathrm{O};\vec{\imath},\vec{\jmath}\)}
\providecommand{\Oijk}{\(\mathrm{O};\vec{\imath},\vec{\jmath},\vec{k}\)}
\providecommand{\Ouv}{\(\mathrm{O};\vec{u},\vec{v}\)}
\providecommand{\intervFF}[2]{\ensuremath{[#1;#2]}}
\renewcommand{\euro}{\ifmmode\text{\texteuro}\else\texteuro\fi}
\DeclareUnicodeCharacter{00B0}{\ifmmode^{\circ}\else\textdegree\fi}
\DeclareUnicodeCharacter{20AC}{\ifmmode\text{\texteuro}\else\texteuro\fi}
% Énoncés : Éducation nationale. Transcription et corrigés : APMEP, Denis Vergès et contributeurs. Découpage et classement : CoopMaths, Sébastien Lozano, Éric Elter et contributeurs. Adaptation de la présentation : Maths Explorer.
% CC BY-SA : https://coopmaths.fr/www/about/
% https://coopmaths.fr/alea/static/bac/2023/tex/bac_2023_03_sujet2_ameriquenord_2.tex
\begin{document}
\noindent\textit{APMEP / CoopMaths. CC BY-SA. Présentation adaptée par Maths Explorer.}\par\medskip
\medskip

On étudie un groupe de \np{3000} sportifs qui pratiquent soit l'athlétisme dans le club A, soit le basketball dans le club B.

En 2023, le club A compte \np{1700} membres et le club B en compte \np{1300}.

\smallskip

On décide de modéliser le nombre de membres du club A et du club B respectivement par deux suites $\left(a_{n}\right)$ et $\left(b_{n}\right)$, où $n$ désigne le rang de l'année à partir de 2023.

L'année 2023 correspond au rang 0. On a alors $a_{0}= \np{1700}$ et $b_{0} = \np{1300}$.

Pour notre étude, on fait les hypothèses suivantes :
\begin{itemize}
\item durant l'étude, aucun sportif ne quitte le groupe ;
\item chaque année, $15\,\%$ des sportifs du club A quittent ce club et adhèrent au club B ;
\item chaque année, $10\,\%$ des sportifs du club B quittent ce club et adhèrent au club $A$.
\end{itemize}

\medskip

\begin{enumerate}
\item Calculer les nombres de membres de chaque club en 2024.
\item Pour tout entier naturel $n$, déterminer une relation liant $a_{n}$ et $b_{n}$.
\item Montrer que la suite $\left(a_{n}\right)$ vérifie la relation suivante pour tout entier naturel $n$, on a : 
\[a_{n+1}= 0,75 a_{n} + 300.\]

\item 
	\begin{enumerate}
		\item Démontrer par récurrence que pour tout entier naturel $n$, on a :
\[1200 \leqslant a_{n+1} \leqslant a_{n} \leqslant \np{1700}.\]
		\item En déduire que la suite $\left(a_{n}\right)$ converge.
	\end{enumerate}
\item Soit $\left(v_{n}\right)$ la suite définie pour tout entier naturel $n$ par $v_{n}=a_{n}- \np{1200}$.
	\begin{enumerate}
		\item Démontrer que la suite $\left(v_{n}\right)$ est géométrique.
		\item Exprimer $v_{n}$ en fonction de $n$.
		\item En déduire que pour tout entier naturel $n, a_{n}= 500 \times 0,75^{n}+ \np{1200}$.
	\end{enumerate}
\item 
	\begin{enumerate}
		\item Déterminer la limite de la suite $\left(a_{n}\right)$.
		\item Interpréter le résultat de la question précédente dans le contexte de l'exercice.
	\end{enumerate}
\item 
	\begin{enumerate}
		\item Recopier et compléter le programme Python ci-dessous afin qu'il renvoie la plus petite valeur de n à partir de laquelle le nombre de membres du club A est strictement inférieur à \np{1280}.

\begin{center}
\psset{unit=1cm}
\begin{tabular}{l l}
\textbf{def}&seuil() :\\
&n = 0\\&A = \np{1700}\\
&\textbf{while ...} :\\
&\qquad n=n+1\\
&\qquad A = \ldots\\
&\textbf{return}\ldots\\
\end{tabular}

\end{center}
		\item Déterminer la valeur renvoyée lorsqu'on appelle la fonction seuil.
	\end{enumerate}
\end{enumerate}

\end{document}
